\documentclass[12pt]{article}

\usepackage[T1]{fontenc}
\usepackage[utf8]{inputenc}
\usepackage{lmodern}
\usepackage[margin=1.25in]{geometry}
\usepackage{setspace}
\usepackage{microtype}
\usepackage{amsmath}
\usepackage{amssymb}
\usepackage{amsthm}
\usepackage[hidelinks]{hyperref}

\theoremstyle{plain}
\newtheorem{theorem}{Theorem}
\newtheorem{lemma}{Lemma}
\newtheorem{proposition}{Proposition}
\newtheorem{corollary}{Corollary}
\theoremstyle{definition}
\newtheorem{definition}{Definition}

\onehalfspacing
\setlength{\parindent}{1.5em}
\setlength{\parskip}{0pt}

\title{The Obligation to Grow}
\author{Flyxion\\Independent Researcher}
\date{November 2026}

\begin{document}
\begin{titlepage}
\centering
\vspace*{1.5cm}

{\LARGE\scshape The Obligation to Grow\par}

\vspace{1cm}

{\large Flyxion\par}
\vspace{0.25cm}
{\normalsize Independent Researcher\par}

\vspace{0.75cm}
{\normalsize November 2026\par}

\vfill

\begin{minipage}{0.88\textwidth}
\begin{center}
{\large\scshape Abstract}
\end{center}

\vspace{0.4cm}

{\sffamily\small
\setstretch{1.15}
\setlength{\parindent}{0pt}

Engagement platforms present the pursuit of an audience as one option among several, yet their defaults, tutorials, milestone notices, and showcase stories treat it as the natural meaning of participation. This essay argues that the shift from option to expectation is best understood at the level of the population rather than the individual. A tactic that improves one account's competitive position can be accurately recommended without that recommendation being justified when everyone adopts it, because prominent visibility is substantially positional, while the aggregate expenditure devoted to securing it can grow without a corresponding growth in audience attention. The essay connects this argument to three earlier ones: one concerning the social pressure that casts every person as a founder, one concerning the archive as a setting in which an artifact can be encountered without invitations to circulate it, and one concerning the advertising-supported platform as the revenue context for the growth prompts described here. The essay closes by proposing a criterion for assessing interfaces: whether they support the activities people came to undertake, as distinct from whether they generate activity. Three formal sections model the generalization argument as a positional adoption game, restate it in categorical and sheaf-theoretic terms, and examine a contest over fixed attention. Three appendices specify archival invariance, derive an externality-correcting charge, and show how repeated solicitation can induce adoption without increasing the user's valuation of the promoted activity.

\par}
\end{minipage}

\vfill
\vspace*{1cm}
\end{titlepage}

\begingroup
\thispagestyle{empty}
\fontfamily{ppl}\selectfont
\vspace*{\fill}
\begin{center}
\begin{minipage}{0.7\textwidth}
\centering
{\Large\itshape Tout le malheur des hommes vient d'une seule chose, qui est de ne savoir pas demeurer en repos dans une chambre.\par}

\vspace{1.2em}

{\large All the unhappiness of men comes from one single thing, which is not knowing how to remain at rest in a room.\par}

\vspace{1.5em}

{\normalsize\scshape Blaise Pascal\par}
\vspace{0.3em}
{\small\itshape Pens\'ees\upshape, Brunschvicg 139 (Lafuma 136)\par}
\vspace{0.2em}
{\footnotesize Translation by the author\par}
\end{minipage}
\end{center}
\vspace*{\fill}
\clearpage
\endgroup

\setcounter{page}{1}

\section{A Reasonable Explanation with an Unanswered Consequence}

A podcast host is deciding whether to publish a conversation he knows to be controversial. He concedes that the evidence offered by his guest is weaker than he would like, he explains why the encounter was nonetheless intellectually interesting, and he admits that the size of his audience influences whom he invites, because the program has to support him. Each of these statements is easy to accept. Hosts of independent programs do have to be paid, a conversation can be worth hearing even when its central claims are poorly supported, and a disclaimer is preferable to leaving disagreement undisclosed, although it is not necessarily preferable to declining publication or substantially revising the presentation. Nothing in the explanation suggests bad faith, and the argument that follows does not depend on suspecting any.

The difficulty appears only when the explanation is multiplied. If a single host adapts his guest list to the demands of attention, his decision is intelligible under the conditions he faces, in the sense that it can be understood from his position and constraints, which is a weaker claim than that it is justified.\footnote{By justified is meant that a decision can be defended after its wider consequences have been weighed. The two are compatible, and the essay's argument turns on keeping them apart.} If a large population of hosts, channel operators, essayists, photographers, and page maintainers all make the same adaptation, the result is an environment in which attracting attention has become the condition under which people can afford to pursue and share their interests at all. The individual explanation remains true in every instance, and it still leaves open what the aggregate is like to live in. The question this essay takes up is therefore not whether a particular host is sincere. It is what social arrangement his explanation makes visible, and what that arrangement asks of people who never intended to enter it.

\section{Participation Acquires an Entrepreneurial Purpose}

Publishing a note, keeping a record, maintaining a friendship, and sharing an enthusiasm need not begin as commercial acts. A great deal of ordinary online activity is of this kind: a person keeps a page because the page is a convenient place to find things again. An interface nevertheless has considerable latitude in how it addresses such a person. It can treat the page as a container for whatever the person puts in it, or it can treat the page as an enterprise at an early stage of development whose natural trajectory is expansion. The earlier essay \emph{Everyone as Founder} \cite{flyxion2026founder} argued that neoliberal societies press individuals and households into founder-like roles, in which survival depends on continual self-promotion and failure is treated as personal inadequacy and not as a systemic shortfall. The present essay follows that pressure into the interface, where it arrives as a purpose supplied to the user before the user has chosen one.

The administrative interface of a Facebook Page offers a concrete instance.\footnote{A Page is the account type that Facebook offers to public figures, businesses, and organizations, as distinct from a personal profile.} A page maintained as a searchable archive of photographs, quotations, and study material is prompted to post on a regular schedule, to respond to other accounts, to thank recent followers by name, to celebrate round-number milestones, and to adopt formats that the platform describes as likely to extend reach. These observations come from the use of a single account over an extended period, and they should be read as such. How widely any particular prompt is deployed, to which accounts, and with what frequency would require documentation that the platform itself holds and that outside observers generally lack. The argument does not need that wider claim. It needs only that the interface, in the cases where it behaves this way, assumes an aspiration the account holder never stated. The archive is treated as an audience-building project awaiting further development, and its modest, finished state is rendered as a shortfall. A purpose supplied in this way must still be sustained, and the interface sustains it by turning the purpose into routine.

\section{The Interface Teaches a Daily Practice}

Growth reaches ordinary routines through a sequence of small obligations rather than through any single demand. Each prompt is manageable on its own terms: a reminder to post, a suggestion to reply to a commenter, a notice that a follower count has crossed a threshold. Their repetition is what matters, because together they establish a continuing discipline of production, measurement, and adjustment. A person who responds to the reminders will, within a short time, be producing on a schedule, reading numerical feedback, and modifying the next piece of work in light of it. That is the daily practice of a small enterprise, acquired without anyone having decided to start one. Milestone notices complete the loop by making numerical expansion feel like progress and by inviting a celebratory post, which is itself new content and a new occasion for measurement.

Beyond the platform's own features, users have built informal institutions around the same expectations. Reciprocal promotion groups,\footnote{Such groups are organized by participants, often in forums or messaging channels adjacent to the platform. The argument concerns their existence as an organizational response to the pursuit of an audience and does not depend on how many people take part.} in which members agree to follow, comment on, share, and sometimes financially support one another's accounts, are a familiar example. They should be distinguished from anything the platform provides, since they are unofficial clubs organized by participants, and their actual economic outcomes are an empirical matter that this essay does not claim to settle. Their relevance lies elsewhere. Audience activity, which might once have been understood as a response that readers happen to give, has become something participants themselves organize and produce, with coordinated schedules and mutual obligations, in anticipation of visibility that may never arrive. Scholars of platform labor have described versions of this development under several names, including free labor in Terranova's sense and the aspirational labor documented by Duffy \cite{terranova2000,duffy2017}. The point to retain here is that an aspiration has been converted into a practice embedded in ordinary routines, and that the practice carries costs in time and attention whether or not the aspiration is ever met.

The revenue structure behind these prompts has been described in \emph{The Hostile Interface} \cite{flyxion2026hostile}, which develops a distinction between a platform in which communication is the service and one in which communication is an input to an advertising business. Where the second arrangement holds, the social graph and the content users produce are valuable chiefly as the material around which advertising is placed, and a user's continued production, together with the attention it attracts, is what the business requires. The growth prompts described above fit this arrangement without any further assumption about intent: a user who is encouraged to post more often, to draw more responses, and to enlarge an audience is supplying what the revenue model consumes, whatever the user's own purposes may be. The present essay takes from that account only the narrower point that, under identifiable competitive and commercial conditions, adaptations that are effective for individuals can sustain outcomes that are burdensome for the group. It does not require that such outcomes be inevitable, that platforms converge on a single form, or that any model of the arrangement identify an optimum.

The arrangement described there has an instructional dimension that Ellul's vocabulary helps to name \cite{ellul1964}. Ellul used the term technique for the totality of rationalized methods optimized for efficiency,\footnote{Ellul's \emph{la technique} is broader than machinery. It names methods in any domain, including administration, advertising, and organization, that are pursued because they are efficient.} and he held that such methods come to organize the surrounding society around their own operational logic, so that people within the system are constrained by that logic and not by their own ends. The interface suggests a refinement. The constraint does not arrive only as a limit on what users can do. It arrives as advice, framed in terms of the user's own aspirations, and it works through encouragement and not prohibition: the prompt to thank new followers or to adopt a format presented as extending reach teaches the participant to reproduce the behavior the business requires, and it is received as help in achieving a goal that may differ from the purposes for which the person first joined.

\section{The Fallacy of Generalizing the Winner's Strategy}

The central analytical claim of the essay concerns a gap between two kinds of advice. Advice about growth can accurately identify what improves one account's competitive position, and it can do so while saying nothing about whether widespread adoption improves the lives of the people who follow it. The statement that short vertical videos increase reach may be true of any given account at a given moment. It does not entail that people ought to make short vertical videos. That recommendation would require an account of what the additional reach enables, what producing it costs, and how widespread adoption changes the conditions under which the advice initially worked. In the same way, the observation that controversial guests attract audiences may be correct for a single program, and it does not establish that the continual escalation of such guests has collective value.

The missing justification is partly a matter of the positional character of prominent visibility.\footnote{A good is positional in Hirsch's sense when its value to the holder depends on how it compares with what others hold. Prominence is an instance, since being seen by many is worth more when others are seen by few.} Attention available to an audience is limited by the hours in a day. Effort is limited too, chiefly by each person's time, but the aggregate expenditure devoted to securing attention can increase without a corresponding increase in the capacity to receive it, because many producers can increase output, refine presentation, and join promotional arrangements while audiences cannot enlarge their hours. When everyone intensifies effort, the rankings that determine who is seen may be left largely unchanged while the cost of holding any given rank rises. This is a familiar structure in economics. Hirsch described goods whose value depends on relative standing as subject to social limits that individual effort cannot overcome, Frank and Cook analyzed markets in which small differences in standing produce large differences in reward and thereby invite excessive investment in competition, and Rosen's earlier work on superstars gives the formal basis for rewards that concentrate at the top \cite{hirsch1976,frankcook1995,rosen1981}. Franck's essay on the economy of attention applies the same logic to attention itself \cite{franck2019}, and Simon's often-cited remark that a wealth of information creates a poverty of attention supplies the underlying scarcity \cite{simon1971}.

An honest assessment of any growth strategy must therefore include what is usually left out: the unsuccessful attempts, the opportunity costs of the hours spent, and the pressure placed on other participants who feel obliged to respond. Commercial failure should be kept distinct from other outcomes, however, since an account can earn nothing while remaining valuable as an archive, a correspondence, or a small community, and the later sections of this essay depend on that distinction. Survivorship shapes the evidence that circulates.\footnote{Survivorship bias is the distortion that arises when the cases that survived some selection process differ systematically from those that did not, and conclusions are drawn only from the survivors.} The accounts that grew are visible and are consulted for lessons, while the larger number that followed similar practices and did not grow are far less likely to be seen, because attention, algorithmic promotion, and the platform's curated showcases all select for success, so that the evidence available to would-be participants is filtered before it reaches them. A recommendation derived from the visible group will tend to overstate the benefit available to the whole. A minimal formal model of this inference is given in Section 8.

The argument should not be pushed into the claim that the system is strictly zero-sum. Useful work does get made, genuine relationships do form, and new audiences do find material they value. The narrower claim is sufficient and easier to defend: individual success cannot by itself justify the general prescription. Whatever the full accounting of benefits turns out to be, the success of the people who prospered is insufficient evidence on its own that the prescription serves the people who were told to follow it. Successful cases can show that something is possible, or can reveal a mechanism, but they cannot establish typical outcomes or justify universal advice.

\section{Exceptional Success Becomes Ordinary Instruction}

The success of those who prospered is insufficient evidence on its own that the prescription serves those who were told to follow it, yet it comes to feel decisive, and the mechanism by which exceptional cases are turned into ordinary instruction deserves examination.

The entrepreneurial invitation is made credible by examples. Videos about rural life, channels devoted to games, expeditions filmed by their participants, and intellectual podcasts all demonstrate that an exceptionally wide range of activity can acquire an audience. This breadth is real, and it is part of what makes the invitation persuasive: nearly any pursuit appears to have a successful exemplar. What the examples cannot show is that the ease of recording and publishing translates into comparable ease of earning. Access to the means of production, in the sense of a camera and a hosting account, has become widely available, although unevenly distributed across differences in equipment, connectivity, free time, accessibility, and personal safety. Access to a sustainable income from using them is not, and the visibility of the former can quietly stand in for the latter.

The successful creator thus has a double role. To viewers, the creator is a source of material. To prospective creators, the same person is a demonstration of an apparently available future. The relevant question is how that second role shapes expectations. The far larger group of unsuccessful participants remains comparatively invisible, producing material that few see and generating no stories that anyone circulates. Platforms are not neutral in this respect, because they choose which stories to feature, and their own self-descriptions emphasize opportunity \cite{gillespie2010,abidin2018}.

A brief comparison with \emph{The Truman Show} illuminates the structure, without carrying the weight of the argument \cite{weir1998}.\footnote{In the 1998 film, the protagonist, Truman Burbank, lives inside a constructed town and does not know that his life is being broadcast to an audience.} The film's protagonist seeks freedom from a life arranged for spectators, and the story treats the observation of his life as the thing from which he must escape. The creator aspiration inverts the direction of the wish: it encourages people to seek recognition by making their own lives watchable, and it frames the arrangement of life for spectators as an achievement. The comparison is useful because it isolates what has changed, namely the valence of being watched, and it should not be taken to suggest that content creators are deceived in the way the fictional character is.

\section{Moral Reservations within a Competitive Medium}

The opening case deserves a second look. The host's disclaimers and methodological objections may express sincere care about the material he is publishing. They can also function as features of the product. An audience that is offered both provocative claims and a visible performance of responsible handling receives more than one that is offered claims alone, and the performance of responsibility can thereby become part of what makes the episode attractive. This observation does not impugn the host's motives, since a gesture can be sincere and also be useful to the one who makes it.

The criticism has to remain specific in order to remain fair. A disclaimer clarifies where the host stands on a claim, and that is a real service to listeners. It does not answer the separate question of whether the selection of the guest, the promotion of the episode, and the framing of the conversation were justified. Likewise, the need to sustain a business explains a choice without settling its wider consequences, because an explanation of why someone acted is a different thing from a defense of what the action contributes to. These are distinct questions, and the argument depends on keeping them apart.

What the case does bring out is a distinction between two kinds of responsibility. One concerns a particular conversation: whether it was handled carefully, whether its limits were disclosed, and whether the audience was treated with respect. The other concerns the reproduction of an environment in which increasingly provocative material becomes an attractive means of survival for anyone who depends on attention. The first can be discharged by a single person acting well. The second cannot, since it belongs to the arrangement as a whole, and it is the arrangement that the interface and its defaults help to constitute.

\section{The Freedom to Leave an Artifact Alone}

The alternative criterion proposed in \emph{Gallery Before Feed} \cite{flyxion2026gallery} can now be stated against this background. A set of photographs can look better in an archive than in a feed, because in the archive each image can be encountered on its own terms, without reaction counters beneath it and without a standing invitation to circulate it. The artifact remains available without continually requesting activity. In a feed, by contrast, the same image is positioned as a bid for response, and its apparent quality becomes entangled with the response it receives. That essay argues that persistent counters, comment prompts, and promotion buttons place evaluation at the moment a viewer should simply recognize the work, and that they turn a timeline from a place of return into a place of obligation. Its alternative treats silence and completion as valid states and not as errors to be corrected.

The same contrast applies to the uses that many people originally had in mind for these services. Retrieval of earlier notes, preservation of a body of work, revisiting an idea after months or years, and spaced repetition of material one wishes to learn are all activities with perfectly good measures of success, and none of those measures is a count of reactions. A record is useful because it can be found and read a long time after it was made. A page can fulfill its purpose completely while its audience remains the same size for a decade. Engagement metrics are poorly placed to register this kind of success, and an interface organized around them will tend to treat it as inactivity.

The normative proposal that follows is modest. People should have the freedom to complete something, to preserve it, and to return to it without being continually recruited into its promotion. Growth is a legitimate goal when a person has a purpose for which growth is a means, such as reaching readers who would benefit from the work, funding further work, or coordinating a community. It becomes an obligation when it is installed as the default meaning of participation, so that the absence of growth must be explained and the absence of promotion must be chosen against a stream of reminders. Growth should require a purpose beyond growth itself.

The design principle that follows is developed in \emph{Gallery Before Feed}, which proposes interaction controls that remain latent until an explicit gesture, a default gallery state without metrics, prompts, or calls to action, and metrics available on demand, with promotion treated as optional (Sections 3, 10.1 to 12.2, and 13.4). Applying that principle to growth prompts is an extension of the earlier essay and not something it separately specifies. Appendix A states the contrast between a quiet default and a quiet invariant, Appendix B gives an externality correction for the adoption model, and Appendix C shows how repeated prompts can make declining costly without changing anyone's preferences.

\section{A Minimal Model of Positional Adoption}

The argument of Section 4 is verbal, and a verbal argument about generalization invites the objection that the conclusion was carried in with the intuition. This section states the simplest model in which the central claim can be checked. It is a model of an inference and not of any platform, and it asserts nothing about the sizes of real prizes or real costs. It also models the generalization argument specifically. It does not model the interface's teaching function, the prompts and routines of Sections 2 and 3, which enter only through the assumption that an account can adopt the practice at a cost.

\begin{definition}[Adoption game]
Let $A=\{1,\dots,n\}$ be a set of $n\ge 2$ accounts and let $p=(p_1,\dots,p_n)$ be a prize vector with $p_1\ge p_2\ge\dots\ge p_n$ that is not constant. Each account $i$ chooses $s_i\in\{0,1\}$, where $s_i=1$ means adopting the growth practice. For a profile $s$ write $S=\{i : s_i=1\}$ and $k=|S|$. Adopters occupy the top $k$ ranks, abstainers occupy the remaining ranks, and the members of each group share the prizes of the ranks they occupy equally.\footnote{Equal sharing within a tied group is an expected-value convention. It gives the same prize as random tie-breaking in expectation.} The positional prize of account $i$ is therefore
\[
\pi_i(s)=
\begin{cases}
T_k,& i\in S,\\
B_k,& i\notin S,
\end{cases}
\qquad
T_k=\frac1k\sum_{j=1}^{k}p_j,\quad B_k=\frac1{n-k}\sum_{j=k+1}^{n}p_j ,
\]
where $T_k$ is defined only for $1\le k\le n$ and $B_k$ only for $0\le k\le n-1$, so that an empty group is never assigned a prize.
Account $i$ has utility $u_i(s)=\pi_i(s)+(v-c)\,s_i$, where $\pi_i(s)$ is its positional prize, $c>0$ is the cost of the practice, and $v\ge 0$ is the value the practice has for the account apart from rank, such as the value of the work as made or of the correspondents it reaches. The constants $c$ and $v$ are the same for all accounts. Write $W(s)=\sum_i u_i(s)$ for total welfare and $m=\frac1n\sum_j p_j$ for the mean prize.
\end{definition}

\begin{lemma}[Conservation of prizes]
For every profile $s$, $\sum_{i}\pi_i(s)=\sum_{j}p_j$. Consequently $W(s)=\sum_j p_j+(v-c)|S|$.
\end{lemma}

\begin{proof}
The $k$ adopters each receive $T_k$, which sum to $\sum_{j\le k}p_j$, and the $n-k$ abstainers each receive $B_k$, which sum to $\sum_{j>k}p_j$. An empty group contributes zero to its sum. The two sums together exhaust the prize vector. Adding the term $(v-c)s_i$ over all $i$ gives the second statement.
\end{proof}

\begin{lemma}[Erosion and the gain from joining]
The sequence $T_1\ge T_2\ge\dots\ge T_n=m$ is nonincreasing, and $T_1>T_n$. For each $k\in\{0,\dots,n-1\}$ the gain from joining when $k$ others have adopted, $D_k=T_{k+1}-B_k$, is strictly positive.
\end{lemma}

\begin{proof}
Since $T_{k+1}=(kT_k+p_{k+1})/(k+1)$ and $p_{k+1}\le p_j$ for every $j\le k$, we have $p_{k+1}\le T_k$, hence $T_{k+1}\le T_k$. The first term is $T_1=p_1$, and $p_1>m$ because $p$ is not constant and $p_1$ is its maximum, so $T_1>T_n$. For the second claim, $T_{k+1}\ge p_{k+1}$ because $T_{k+1}$ averages entries that are all at least $p_{k+1}$, and $B_k\le p_{k+1}$ because $B_k$ averages entries that are all at most $p_{k+1}$. Equality in both would force $p_1=\dots=p_{k+1}=\dots=p_n$, which is excluded. Therefore $D_k>0$.
\end{proof}

Let $\delta=\min_k D_k$, which is positive because the minimum is over finitely many positive numbers.

\begin{theorem}[Individually correct advice, collectively worse outcome]
Suppose $c-v<\delta$. Then adopting is a strictly dominant strategy for every account, and the profile in which all accounts adopt is the unique Nash equilibrium.\footnote{A profile is a Nash equilibrium if no account can raise its own utility by changing only its own choice.} If in addition $v<c$, then every account has strictly lower utility when all adopt than when none adopts, and total welfare falls by $n(c-v)$.
\end{theorem}

\begin{proof}
Let $i\notin S$ with $|S|=k$. Joining changes the utility of $i$ from $B_k$ to $T_{k+1}+(v-c)$, a change of $D_k+v-c\ge \delta+v-c>0$. This holds for every $S$ and every $i$, so adopting is strictly dominant. In any Nash equilibrium every account plays its strictly dominant strategy, so the equilibrium is unique and equals the all-adopt profile. At that profile each account receives $T_n+v-c=m+v-c$, while at the profile in which none adopts each receives $B_0=m$. If $v<c$ the former is strictly smaller for every account, and by Lemma 1 the total changes by $n(v-c)$.
\end{proof}

\begin{corollary}[A purpose beyond growth]
For every prize vector $p$ and every profile, the welfare effect of one further adoption is exactly $v-c$. Redistribution of the fixed positional prizes contributes nothing to the change in total welfare, so the only term that can make adoption collectively worthwhile is the non-positional value $v$ exceeding the cost $c$.
\end{corollary}

The corollary is the formal counterpart of the sentence that growth should require a purpose beyond growth itself. It also explains why the advice in Theorem 1 is not mistaken in any individual case: the recommendation to adopt is correct for each account whatever the others do, and its generalization is the thing that fails. The model leaves several features fixed, namely binary effort, identical accounts, a fixed prize vector, and no entry of new audience. The assumption of a fixed prize vector is the one that matters most. Additional reach for one account does not by itself enlarge the aggregate prize, since it can simply redistribute existing attention. Conservation in Lemma 1 would fail if adoption changed the total prize pool, for example by increasing total monetizable attention. Additional non-positional benefits not already represented by $v$ would instead change the welfare calculation without affecting the lemma. Whether such changes offset the cost of adoption is an empirical question, and it is the one that Section 4 places at the center of any recommendation to produce more.

\section{Two Structural Reformulations}

The model of Section 8 can be restated in two further vocabularies, which make its structure explicit without adding assumptions. The first uses the category of adopter sets, and the second uses sheaves over the set of accounts.

\subsection{The Hypercube of Adoption}

The profiles of the previous section form a familiar category. Let $\mathbf{B}$ denote the poset $(\mathcal P(A),\subseteq)$ regarded as a category, whose objects are sets of adopters and whose morphisms are inclusions \cite{maclane1998}. Its generating morphisms are the edges $S\to S\cup\{i\}$ with $i\notin S$, and these are the edges of the $n$-dimensional hypercube. Three functions on those edges describe what an adoption does.

\begin{definition}
For an edge $(S,i)$ with $i\notin S$ let
\begin{align*}
g(S,i)&=u_i(S\cup\{i\})-u_i(S),\\
w(S,i)&=W(S\cup\{i\})-W(S),\\
\varepsilon(S,i)&=w(S,i)-g(S,i).
\end{align*}
These are the individual gain, the social gain, and the externality, respectively. By definition $\varepsilon(S,i)=\sum_{j\ne i}\bigl(u_j(S\cup\{i\})-u_j(S)\bigr)$, the total change in the utilities of everyone other than the adopter.
\end{definition}

\begin{proposition}
For every edge $(S,i)$ with $|S|=k$, the social gain is $w(S,i)=v-c$, the individual gain is $g(S,i)=D_k+v-c$, and the externality is $\varepsilon(S,i)=-D_k$. Moreover $w$ is the coboundary of the function $W$ on vertices, so for every maximal chain $\emptyset=S_0\subset S_1\subset\dots\subset S_n=A$ the sum of social gains along the chain equals $W(A)-W(\emptyset)=n(v-c)$, whichever order the accounts adopt in.
\end{proposition}

\begin{proof}
The first statement is Lemma 1, since $W(S\cup\{i\})-W(S)=(v-c)$. The second repeats the computation in the proof of Theorem 1. The third follows by subtraction. The last statement holds because $w(S,i)=W(S\cup\{i\})-W(S)$ is a difference of vertex values, so the sum along any chain telescopes.
\end{proof}

The third formula is the substance of the proposition. The externality that an adoption imposes on everyone else is exactly the negative of the positional gain the adopter obtains, so the whole of that gain is a transfer and none of it is a creation of value. This is the precise sense in which the visibility won by the account that adopts is lost by the accounts that do not. It is also the answer to the objection that Section 4 declined to treat the system as zero-sum: the model is zero-sum in its positional part and not in the part measured by $v-c$.

\begin{corollary}
If $c-v<\delta$ and $v<c$, then $g>0$ and $w<0$ on every edge of the hypercube. Every maximal chain of individually advantageous adoptions therefore ends in a state worse for every account than its starting point, and the total social change along the chain is independent of the order of adoption.
\end{corollary}

This section is a structural reformulation of the model in Section 8 using the Boolean lattice of adopter sets. It makes the externality and the path independence of social gain explicit, and it is not an independent explanation of why interfaces induce adoption.

\subsection{Sections, Gluing, and the Locality of Payoff}

A second reformulation addresses why a recommendation that is correct and verifiable in each case can fail as a policy, and sheaf theory gives a precise vocabulary for it, with a modest result \cite{maclanemoerdijk1992}. Let $X=A$ carry the discrete topology,\footnote{A sheaf assigns data to regions of a space so that data on overlapping regions agree and can be glued into data on their union. On a finite discrete space the regions are just sets of points, and a sheaf reduces to a family of choices indexed by the points.} and define the sheaf of verdicts $F$ by $F(U)=\prod_{i\in U}\{0,1\}$, with restriction maps given by projection. A local section over $\{i\}$ is one account's decision or the recommendation addressed to it, and a global section is a profile. Because $X$ is discrete, a sheaf on $X$ is the same as a family of sets indexed by the points, equivalently a functor on the discrete category on $A$, and its global sections form the limit of that diagram, which is the product.

\begin{proposition}[Gluing is unobstructed]
Every family of local sections $(s_i)_{i\in A}$ with $s_i\in F(\{i\})$ glues to a unique global section $s\in F(A)$. For any sheaf of abelian groups on the finite discrete space $X$, the higher sheaf cohomology groups vanish.
\end{proposition}

\begin{proof}
The singletons are pairwise disjoint, so the compatibility condition on overlaps is vacuous, and the sheaf axiom says that $F(A)\to\prod_i F(\{i\})$ is a bijection. For the second statement, the global sections functor on a finite discrete space is a finite product of stalk functors, hence exact, and an exact global sections functor has vanishing higher derived functors.
\end{proof}

The proposition says that local advice always assembles into a population-wide practice. Nothing in the cohomology of the space obstructs the assembly, and the failure described in the previous sections is not of that kind. It lies in the evaluation of a section, which the next proposition makes exact.

\begin{proposition}[Payoff is not local]
Let $\rho_i:F(A)\to F(\{i\})$ be the restriction to the point $i$. The utility $u_i:F(A)\to\mathbb R$ does not factor through $\rho_i$. In particular, for every $i$,
\[
u_i(\{i\})-u_i(A)=p_1-m>0 ,
\]
where the arguments denote the profiles in which only $i$ adopts and in which every account adopts.
\end{proposition}

\begin{proof}
Both profiles restrict to the same section $s_i=1$ at $i$. The utility of $i$ is $T_1+(v-c)=p_1+(v-c)$ in the first and $T_n+(v-c)=m+(v-c)$ in the second, so the difference is $p_1-m$, which is positive by Lemma 2. Two global sections with the same restriction to $\{i\}$ therefore give $i$ different utilities.
\end{proof}

The proposition gives a precise form, in the language of sections, to the dependence that Section 4 describes. Payoff dependence alone does not imply that generalization is harmful, and the welfare comparison of Theorem 1 supplies that further conclusion. A recommendation is a local section, and the evidence for it was gathered at a particular global section, namely the environment in which the exemplary account succeeded. The value of the local section is not determined by itself. It depends on the global section into which it is embedded, and Lemma 2 shows the direction of the dependence: the prize of an adopter is nonincreasing in the number of adopters and strictly smaller at full adoption than for a lone adopter. Gluing guarantees that the population can adopt the advice. It does not guarantee that the value observed at one global section survives at the glued one, and the premium enjoyed by a lone adopter is exactly what full adoption removes.

\section{A Contest Model with Attention Shares}

The rank-based model of Section 8 treats prominence as a rank. A second stylized model treats it as a share of a fixed quantity of attention,\footnote{The share function is of the ratio form often called a Tullock contest. Each account's share is its own effort divided by total effort.} which allows effort to vary continuously and gives a quantitative statement of the dissipation involved. As before, it models an inference and does not describe any platform.

\begin{definition}[Attention contest]
Let $N\ge2$ accounts choose efforts $e_i\ge0$. Total audience attention is a fixed quantity $K$, and account $i$ receives the share $R_i(e)=e_i/\sum_{j}e_j$, so that its attention is $a_i=KR_i$ and $\sum_i a_i=K$ whatever the efforts are; at the all-zero profile set $R_i=1/N$. The share is worth $V>0$ per unit of share to the account, and effort costs $c>0$ per unit, so that the payoff is $\pi_i(e)=VR_i(e)-ce_i$.
\end{definition}

\begin{theorem}[Symmetric equilibrium]
The profile in which every account chooses $e^*=\dfrac{V(N-1)}{cN^2}$ is a Nash equilibrium. At that profile each account receives the share $1/N$ and the payoff $V/N^2$, aggregate effort is $E^*=\dfrac{V(N-1)}{cN}$, aggregate cost is $cE^*=V\dfrac{N-1}{N}$, and the total payoff is $V-cE^*=V/N$.
\end{theorem}

\begin{proof}
Fix $i$ and let $Y=\sum_{j\ne i}e_j=(N-1)e^*>0$. The payoff of $i$ as a function of its own effort is $f(e)=Ve/(e+Y)-ce$, whose second derivative $-2VY/(e+Y)^3$ is negative, so $f$ is strictly concave on $e\ge0$. Its derivative at $e^*$ is $VY/(e^*+Y)^2-c=V(N-1)e^*/(Ne^*)^2-c=V(N-1)/(N^2e^*)-c=0$. A stationary point of a strictly concave function on $[0,\infty)$ is its global maximizer, so no deviation of $i$ is profitable. The remaining quantities follow by substitution: $\pi_i=V/N-ce^*=V/N^2$, and $\sum_i VR_i=V$.
\end{proof}

\begin{corollary}[Dissipation under a fixed budget]
The share of the contested value consumed by effort is $cE^*/V=(N-1)/N$, which tends to $1$ as $N$ grows, while total attention remains $K$ for every effort profile. Privately rational competition can therefore dissipate most of the value being contested without enlarging the attention being contested.
\end{corollary}

Two cautions govern the reading of this result. The first concerns the exponent in the common generalization $R_i=e_i^{\alpha}/\sum_j e_j^{\alpha}$. The symmetric stationary candidate is then $e^*=V\alpha(N-1)/(cN^2)$, with payoff $V\bigl(N-\alpha(N-1)\bigr)/N^2$, and for $\alpha>N/(N-1)$ this payoff is negative, whereas zero effort against positive opponents yields zero. The candidate then cannot be an equilibrium, so the baseline above sets $\alpha=1$ and does not rely on increasing returns to effort. The second caution concerns the cost parameter. As $c\to0^+$ the effort units $E^*$ diverge, but the total cost $cE^*=V(N-1)/N$ does not. If automation makes each unit of effort cheaper, more units need not mean more exhaustion or energy use, since those consequences would need their own terms for time, energy, or opportunity cost.

The conclusion is correspondingly narrow. The model does not show that the expenditure is pure waste, because it assigns no value to the content, learning, or relationships that the effort may produce, and Section 4 already allows for such benefits. It shows that, once attention is a fixed budget, the contest component of the activity returns little to the participants as a group, and that the justification for the effort must be found in what lies outside the contest.

\section{Conclusion: An Account Need Not Become an Enterprise}

The essay has distinguished two things that are easily run together: the availability of an entrepreneurial path and its installation as the default meaning of participation. The first is a good. People should be able to pursue audiences and commercial success, and a platform that makes this possible provides something of value. The second is a different matter, since it converts an option into an expectation and thereby changes what it costs to decline. Declining ceases to be the ordinary state of having a page and becomes a visible failure to take advantage of the tools provided.

The generalization argument explains why this conversion deserves scrutiny and not simple deference to the success stories. A recommendation that is correct for each account taken singly can be mistaken for the population, because competition for prominent visibility can increase aggregate effort without producing a corresponding increase in available audience attention. The host in the opening case made a choice that is intelligible given the conditions he faced, and the same is true of most of the people who follow a platform's advice, though intelligibility is not the same as justification. The conditions are what remain to be assessed.

The criterion that results is simple to state and demanding to apply. Measuring how much activity a platform produces cannot establish whether it serves the purposes for which people arrived, and the two should not be allowed to stand in for one another. Archiving, correspondence, study, and the keeping of a body of work are among those activities, and for them the relevant questions are whether the thing made is available later, whether it can be found, and whether its maker can leave it alone. An interface can be assessed against these questions, and so can the instructions it gives its users.

The earlier essays supply complementary parts of this conclusion. \emph{Everyone as Founder} describes how social and economic pressure casts individuals and households as ventures that must continually promote themselves, \emph{Gallery Before Feed} describes what is preserved when an artifact is left alone, and \emph{The Hostile Interface} describes the revenue structure within which the prompts arrive. The present essay adds the population-level argument that supplies the justification for treating the choice as the user's and not the platform's. The formal sections state that argument in three results. The welfare effect of any adoption equals its non-positional value minus its cost, so that only a purpose beyond growth can make it collectively worthwhile. Under a fixed attention budget the contest component of the effort can consume most of what is contested without enlarging it. And repeated prompts can raise the cost of declining enough that observed adoption need not reveal that the user values the promoted activity. Because the success of some cannot justify the prescription for all, growth cannot be assumed to be the right objective for everyone; and because users differ in their purposes, and those purposes should govern whether growth is an appropriate objective for them, an interface should preserve the user's choice between the two postures instead of installing one as the default. At the level of design, this means that growth tools remain available by deliberate choice, while archival use requires neither promotion nor an explanation for declining it. An account need not become an enterprise, and a service that treats every account as one has answered a question that its users did not ask.

\appendix

The three appendices extend the formal sections. Their results concern specified models, and their application to an actual platform would require evidence.

\section{Archival Autonomy as an Interface Invariant}

The freedom to maintain an archive can be specified as a property of permitted interface transitions, and not only as the availability of a preference setting.

Let $Z$ be a set of interface states, $U$ a set of user actions, and $P$ a set of autonomous platform actions, including scheduled notifications and unsolicited recommendations.\footnote{The symbol $\sqcup$ denotes disjoint union, so that no action is both a user action and a platform action.} Let
\[
\tau:Z\times(U\sqcup P)\longrightarrow Z
\]
be the transition function. Let $Q\subseteq Z$ denote the quiet archival states, and let $U_0\subseteq U$ contain ordinary archival actions: reading, searching, navigation, and preservation. Let $U_1\subseteq U$ contain explicit requests to enter a growth-oriented mode, and assume $U_0\cap U_1=\varnothing$. The set $U_0$ names the ordinary archival actions for reference, while the formal condition below quantifies over all actions outside $U_1$.

\begin{definition}[Archival invariance]
The interface preserves archival autonomy if
\[
\tau\bigl(Q,(U\setminus U_1)\sqcup P\bigr)\subseteq Q.
\]
Departure from the quiet region therefore requires an explicit request in $U_1$. Since $U_0\cap U_1=\varnothing$, the condition covers ordinary archival actions and autonomous platform activity, and it also excludes any other user action outside $U_1$ from triggering growth mode by accident.
\end{definition}

\begin{proposition}[Persistence without renewed refusal]
If archival invariance holds and $z_0\in Q$, then every finite execution from $z_0$ containing no action in $U_1$ remains in $Q$.
\end{proposition}

\begin{proof}
Induct on the length of the execution. The initial state belongs to $Q$. If the current state belongs to $Q$, archival invariance places the next state in $Q$, and the conclusion follows.
\end{proof}

The distinction between a quiet default and an invariant is substantial. A default describes the initial state, whereas an invariant constrains what the platform may subsequently do without renewed authorization. Let $h:Z\to\{0,1\}$ indicate whether unsolicited growth solicitation is visible, and require $h|_Q=0$. If a scheduled platform action $p\in P$ takes a state $q\in Q$ to a state with $h(\tau(q,p))=1$, the interface fails archival invariance even if the user can dismiss the solicitation immediately, because the resulting state lies outside $Q$. This supplies a precise design criterion: remaining quiet must not require a recurring act of refusal.

\section{Externality-Corrected Growth Advice}

The adoption model of Section 8 permits an exact intervention. Rather than asking whether a growth practice benefits its adopter, an interface could evaluate the change in welfare that the adoption produces. Retain the notation of Section 8, and suppose that an adopter pays a charge $t(k)$ when the final number of adopters is $k$. Charges are transferred to an external public account, and participants do not receive those transfers within this model. Set $t(0)=0$ by convention, so that the charge is defined at $S=\varnothing$, where no account pays it. Define adjusted utility by
\[
\widetilde u_i(S)=u_i(S)-\mathbf{1}_{\{i\in S\}}\,t(|S|).
\]

\begin{proposition}[Exact correction]
Set $t(k)=D_{k-1}$ for $1\le k\le n$. Then the adjusted gain from adopting, for every account and every configuration of the other accounts, is exactly $v-c$.
\end{proposition}

\begin{proof}
Let $k$ other accounts have adopted. Abstention incurs no charge. Adoption changes the original utility by $D_k+v-c$ and incurs the charge $t(k+1)=D_k$, so the adjusted difference is $v-c$.
\end{proof}

\begin{corollary}
If $v<c$, abstention is strictly dominant in the corrected game. If $v>c$, adoption is strictly dominant. If $v=c$, each account is indifferent between its two actions.
\end{corollary}

The charge equals the magnitude of the negative marginal externality imposed on the other accounts,
\[
t(k)=D_{k-1}=-\varepsilon(S,i),\qquad |S|=k-1,
\]
where $\varepsilon$ is the quantity computed in Section 9. The distinction matters because the externality is negative and the corrective charge is positive, and the exactness of the correction depends on the identity holding at every $k$. The correction aligns each account's private gain with the change in the original welfare function $W$, counting the public receipts as transfers. It does not establish that the sum of the participants' adjusted utilities increases. It removes the incentive to adopt solely for a positional transfer while retaining the incentive to create non-positional value. This is an implementation result under strong informational assumptions. It presupposes knowledge of the prize structure and excludes strategic effects from redistributing the charges to participants, and it establishes neither the desirability of an actual tax nor a practical method of estimating one. Its conceptual consequence is nonetheless useful: growth advice that measures only the adopter's gain omits exactly the term required to align the private decision with the social change.

\section{Repeated Solicitation as a Cost of Remaining Still}

Persistent prompts can alter participation without changing the value a user assigns to growth. Consider an account holder who prefers quiet archival use. While the account remains quiet, growth prompts arrive according to a Poisson process of rate $\lambda>0$,\footnote{A Poisson process generates events at random times, with $\lambda$ events per unit time on average, and events in disjoint time intervals are independent.} each prompt imposes a dismissal or interruption cost $d>0$, and future costs are discounted at rate $r>0$. Suppose that adopting the promoted mode stops these particular prompts but imposes a one-time net loss $L>0$, incorporating the user's evaluation of the change and its subsequent consequences. Preferences, $\lambda$, and $d$ remain fixed.

\begin{proposition}[The price of continued refusal]
The expected discounted cost of remaining quiet indefinitely is
\[
C_Q=\frac{\lambda d}{r}.
\]
Consequently, whenever $0<L<\lambda d/r$, immediate adoption is preferable to indefinite refusal, although adoption would be undesirable in the absence of solicitation.
\end{proposition}

\begin{proof}
Let $N_t$ count prompts up to time $t$, a process with intensity $\lambda$. Then
\[
C_Q=\mathbb{E}\left[\int_0^\infty e^{-rt}\,d\,\mathrm{d}N_t\right]
=\lambda d\int_0^\infty e^{-rt}\,\mathrm{d}t=\frac{\lambda d}{r}.
\]
Immediate adoption costs $L$, so the stated inequality gives the comparison.
\end{proof}

The result requires no persuasion, no misinformation, and no prohibition on archival use. The interface changes the cost of declining while leaving the user's underlying valuation of adoption negative. Under the additional restriction that adoption occurs at a deterministic time $\theta$, the expected discounted cost is
\[
C(\theta)=\frac{\lambda d}{r}\bigl(1-e^{-r\theta}\bigr)+e^{-r\theta}L,
\qquad
C'(\theta)=e^{-r\theta}(\lambda d-rL).
\]
If $\lambda d>rL$, the cost is increasing in $\theta$ and immediate adoption minimizes it, and if $\lambda d<rL$, the cost is decreasing and indefinite refusal minimizes it. The threshold depends on the architecture of solicitation and not on any improvement in the promoted activity.

The assumptions matter. Actual prompts may be harmless, may remain after adoption, or may be permanently disabled, and where those alternatives are available without comparable cost they change the decision problem. The model identifies a mechanism to test: repeated solicitation can make an unwanted transformation preferable to the burden of continually declining it, so that observed adoption need not reveal that the user values the promoted activity.

\begin{thebibliography}{99}

\bibitem{abidin2018}
Abidin, Crystal.
\textit{Internet Celebrity: Understanding Fame Online}.
Bingley: Emerald Publishing, 2018.

\bibitem{duffy2017}
Duffy, Brooke Erin.
\textit{(Not) Getting Paid to Do What You Love: Gender, Social Media, and Aspirational Work}.
New Haven: Yale University Press, 2017.

\bibitem{ellul1964}
Ellul, Jacques.
\textit{The Technological Society}.
Translated by John Wilkinson.
New York: Alfred A. Knopf, 1964.

\bibitem{flyxion2026founder}
Flyxion.
\textit{Everyone as Founder: Colonial Roots, Patriarchy, and Field-Theoretic Society}.
Essay, September 7, 2025.
\url{https://standardgalactic.github.io/gamma-waves/economics/essay.pdf}.

\bibitem{flyxion2026gallery}
Flyxion.
\textit{Gallery Before Feed: Why Engagement-Optimized Interfaces Destroy Coherence and How Merge-Aware Systems Restore Meaning}.
Essay, January 10, 2026.
\url{https://standardgalactic.github.io/antivenom/opinion/Gallery\%20Before\%20Feed.pdf}.

\bibitem{flyxion2026hostile}
Flyxion.
\textit{The Hostile Interface: Advertising Saturation and the Collapse of the Social Information Environment}.
Essay, January 19, 2026.
\url{https://standardgalactic.github.io/antivenom/The\%20Hostile\%20Interface.pdf}.

\bibitem{franck2019}
Franck, Georg.
The Economy of Attention.
\textit{Journal of Sociology} 55, no. 1 (2019): 8--19.
Originally published in German in \textit{Merkur}, 1993.

\bibitem{frankcook1995}
Frank, Robert H., and Philip J. Cook.
\textit{The Winner-Take-All Society}.
New York: Free Press, 1995.

\bibitem{gillespie2010}
Gillespie, Tarleton.
The Politics of Platforms.
\textit{New Media \& Society} 12, no. 3 (2010): 347--364.

\bibitem{hirsch1976}
Hirsch, Fred.
\textit{Social Limits to Growth}.
Cambridge, MA: Harvard University Press, 1976.

\bibitem{maclane1998}
Mac Lane, Saunders.
\textit{Categories for the Working Mathematician}.
2nd ed.
New York: Springer, 1998.

\bibitem{maclanemoerdijk1992}
Mac Lane, Saunders, and Ieke Moerdijk.
\textit{Sheaves in Geometry and Logic: A First Introduction to Topos Theory}.
New York: Springer, 1992.

\bibitem{rosen1981}
Rosen, Sherwin.
The Economics of Superstars.
\textit{American Economic Review} 71, no. 5 (1981): 845--858.

\bibitem{simon1971}
Simon, Herbert A.
Designing Organizations for an Information-Rich World.
In Martin Greenberger, ed.,
\textit{Computers, Communications, and the Public Interest},
37--72.
Baltimore: Johns Hopkins Press, 1971.

\bibitem{terranova2000}
Terranova, Tiziana.
Free Labor: Producing Culture for the Digital Economy.
\textit{Social Text} 18, no. 2 (2000): 33--58.

\bibitem{weir1998}
Weir, Peter, dir.
\textit{The Truman Show}.
Paramount Pictures, 1998.

\end{thebibliography}

\end{document}
